\section{数据结构}

\subsection{\texttt{page\_directory\_t} 与 \texttt{page\_table\_t}}

\codefile{src/include/kernel/vmm.h}
\begin{lstlisting}[style=cstyle]
typedef struct {
    uint32_t entries[1024];  /* 1024 个 PDE，每个 4 字节 */
} __attribute__((aligned(4096))) page_directory_t;

typedef struct {
    uint32_t entries[1024];  /* 1024 个 PTE，每个 4 字节 */
} __attribute__((aligned(4096))) page_table_t;
\end{lstlisting}

\begin{figure}[H]
  \centering
  % figures/ch08/fig04-pd-layout.tex — auto-extracted from chapter source
\begin{tikzpicture}[
  cell/.style={draw, minimum width=5cm, minimum height=0.5cm, font=\ttfamily\small, anchor=north},
  addr/.style={font=\ttfamily\scriptsize\color{gray}, anchor=east},
]
  \node[font=\bfseries] at (2.5, 1) {page\_directory\_t (4\,KiB)};
  \node[cell, fill=red!10]  (e0) at (0, 0.3) {entries[0] — PDE 0};
  \node[cell, fill=red!10]  (e1) at (0, -0.2) {entries[1] — PDE 1};
  \node[cell, fill=red!10]  (e2) at (0, -0.7) {entries[2] — PDE 2};
  \node[font=\small] at (2.5, -1.2) {$\vdots$};
  \node[cell, fill=yellow!15] (e768) at (0, -1.7) {entries[768] — \hex{C0000000}};
  \node[font=\small] at (2.5, -2.2) {$\vdots$};
  \node[cell, fill=red!10]  (e1023) at (0, -2.7) {entries[1023] — PDE 1023};

  % 标注
  \node[addr] at (-0.1, 0.05) {+\hex{000}};
  \node[addr] at (-0.1, -0.45) {+\hex{004}};
  \node[addr] at (-0.1, -0.95) {+\hex{008}};
  \node[addr] at (-0.1, -1.95) {+\hex{C00}};
  \node[addr] at (-0.1, -2.95) {+\hex{FFC}};

  % 花括号标注 identity / high-half
  \draw[decorate, decoration={brace, amplitude=5pt}] (5.1, 0.3) -- (5.1, -0.7)
    node[midway, right=6pt, font=\scriptsize] {Identity map};
  \draw[decorate, decoration={brace, amplitude=5pt}] (5.1, -1.7) -- (5.1, -2.7)
    node[midway, right=6pt, font=\scriptsize] {High-half (内核)};
\end{tikzpicture}

  \caption{\texttt{page\_directory\_t} 内存布局——PDE[768] 对应 \hex{C0000000}}
  \label{fig:pd-layout}
\end{figure}

\begin{whybox}
为什么 \texttt{page\_directory\_t} 用 \texttt{\_\_attribute\_\_((aligned(4096)))}？

\reg{CR3} 要求页目录物理地址必须 4\,KiB 对齐（低 12 位为 0）。
用 \texttt{aligned} 属性声明的静态变量由链接器保证对齐。
静态变量放在 BSS 段，启动时自动清零——所有 PDE 初始 Present=0，安全。
\end{whybox}

\begin{thinkbox}
每个 PDE 覆盖 $1024 \times 4\,\text{KiB} = 4\,\text{MiB}$ 的虚拟空间。
PDE[768] 对应虚拟地址 $768 \times 4\,\text{MiB} = \hex{C0000000}$
——正好是 High-Half Kernel 的起始地址。这是巧合吗？
（答：不是。\hex{C0000000} 被选为内核基地址，正因为它恰好对齐到 PDE 边界，
简化了地址空间划分。）
\end{thinkbox}

% ============================================================================

